% Timing: 13:00--25:00   (paper: Sec. 1.3)
\section{Technical overview}

% ============================================================
\begin{frame}{Proof outline (at high-level)}

\KEY{Goal Lemma} \ Every \CoR{} term equivalent to $G_k$ needs $2^{\Omega(k)}$ compositions.

\vfill
We construct structures $\M_0$ and $\M_1, \dots, \M_{2^k}$ such that
\begin{itemize}
\uncover<2->{%
  \item $G_k$ distinguishes all of $\M_1, \dots, \M_{2^k}$ from $\M_0$:
    \begin{center}
    \renewcommand{\arraystretch}{1.15}
    \begin{tabular}{@{}w{l}{.6cm}|w{c}{.78cm}|*{7}{w{c}{.78cm}}|w{c}{.9cm}|w{c}{1.25cm}@{}}
      & $\M_0$ & $\M_{1}$ & $\M_{2}$ & \multicolumn{4}{w{c}{\dimexpr 3.12cm+6\tabcolsep\relax}}{$\cdots$} & $\M_{2^k}$ & \#diff. & \uncover<6->{\#comp.\,$\ge$} \\
      \hline
      $G_k$ & T & \alert{F} & \alert{F} & \multicolumn{4}{w{c}{\dimexpr 3.12cm+6\tabcolsep\relax}}{$\cdots$} & \alert{F} & \alert{$2^k$} & \uncover<6->{$2^k/2$} \\
    \end{tabular}
    \end{center}
}
\uncover<3->{%
   \item 
\KEY{Key Lemma} \ If a \CoR{} term $t$ distinguishes $d$ structures among $\M_1, \dots, \M_{2^k}$ from $\M_0$,
then $t$ needs $\ge d/2$ compositions.
\uncover<4->{%
\begin{center}
\renewcommand{\arraystretch}{1.15}
\begin{tabular}{@{}w{l}{.6cm}|w{c}{.78cm}|*{7}{w{c}{.78cm}}|w{c}{.9cm}|w{c}{1.25cm}@{}}
  & $\M_0$ & $\M_{1}$ & $\M_{2}$ & $\M_{3}$ & $\M_{4}$ & $\M_{5}$ & $\cdots$ & $\M_{2^k}$ & \#diff. & \#comp.\,$\ge$ \\
  \hline
  $t$ & T & \alert{F} & T & \alert{F} & T & T & $\cdots$ & T & \alert{$2$} & $1$ \\
  $t'$ & F & \alert{T} & \alert{T} & \alert{T} & \alert{T} & F & $\cdots$ & F & \alert{$4$} & $2$ \\
\end{tabular}
\end{center}
}}
\end{itemize}

\vfill
\uncover<6->{\HP{if $t \equiv G_k$, then $t$ needs $\ge 2^k/2$ compositions.}}

\vfill
\uncover<7->{\large \KEY{TASK} How can we construct such $\M_0$ and $\M_1, \dots, \M_{2^k}$?}

\note{
\textbf{スライド概要}：$G_k$ が $\M_0$ で成り立ち、$\M_1, \dots, \M_{2^k}$ で成り立たない構造を作る。$d$ 個を $\M_0$ から見分ける項は、合成を $d/2$ 個以上必要とする。

\textbf{メモ}：
\begin{itemize}
  \item 表の T/F は、各 $\M_i$ に対応する点対での値（論文の test pair）
  \item 以降：$\M_0$ の作り方、$\M_i$ ＝ hub での塗り替え $\M_\lambda$、合成が新たに見分けるのは高々 2 個である理由
\end{itemize}
}
\end{frame}

% ============================================================
% The structure M_0 (edge-colored complete graph, two middle vertices, the hub)
\begin{frame}{Structures as colored complete graphs}

\KEY{View} structures as \KEY{complete graphs}, colored with $2^k$ colors $\mu \subseteq [k]$:
$P_i \iff i \in \mu$.
\\ \NOTE{$Q_i$ is used as the complement of $P_i$ (i.e., $i \notin \mu$): exactly one of $P_i, Q_i$ holds on each edge.}

\vfill
\begin{center}
\begin{tikzpicture}[every node/.style={font=\small}]
  % same style as the figure of G_k (example with k = 3): faint candidates for z with dashed edges
  % (the dashed edges are drawn only near the candidates, to leave room for the label of x-z)
  \node[pv,minimum size=7pt,label=left:$x$] (x) at (0,0) {};
  \node[pv,minimum size=7pt,label=right:$y$] (y) at (4.8,0) {};
  \foreach \h in {2.45,3.0} {
    \node[pv,minimum size=5pt,draw=black!30] (c) at (3.0,\h) {};
    \draw[densely dashed,black!30,shorten <=2.2cm] (x) -- (c);
  }
  \node[font=\footnotesize,text=black!60] at (3.45,2.72) {$\exists z$};
  \node[pv,minimum size=7pt,label=right:$z$] (z) at (3.0,1.9) {};
  \draw[thick,cA] (x) -- node[below,text=cA]
    {color $\mu = \set{1, 3}$: \ $P_1, \ Q_2, \ P_3$} (y);
  \draw[thick,cB] (x) -- node[pos=.45,above left,align=right,text=cB,font=\footnotesize]
    {color $\mu^c = \set{2}$:\\ $Q_1, \ P_2, \ Q_3$} (z);
  \node[anchor=west,align=left] at (5.8,1.0)
    {for every $i \in [3]$,\\
    the $x$-$y$ and $x$-$z$ edges\\
    have labels $P_i$ and $Q_i$, in either order\\
    $\Rightarrow$ \ $G_3(x, y)$ \ \NOTE{($k = 3$)}};
\end{tikzpicture}
\end{center}

\vfill
\HP{$G_k(x, y)$ holds iff some $x$-$z$ edge has the color $\mu^c$, where $\mu$ is the color of the $x$-$y$ edge.}

\note{
\textbf{スライド概要}：構造を色付き完全グラフとして見る。$G_k$ は「補の色の $x$-$z$ 辺がある」と読める。

\textbf{メモ}：
\begin{itemize}
  \item 各辺で $P_i$ と $Q_i$ のちょうど一方が成り立つ
  \item 例：$k = 3$、$\set{1,3}$ と $\set{2}$ は互いに補
\end{itemize}
}
\end{frame}

% ============================================================
\begin{frame}{The randomly colored structure $\M_0$}

\begin{columns}[c]
  \begin{column}{.36\textwidth}
    \centering
    % Overall shape of M_0 (used on the slide 'The randomly colored structure M_0').
% Edge colors are fixed pseudo-random choices (seed 7); u (left), v (right) and the
% two middle vertices z (upper), z' (lower) of the pattern u -mu- z -nu- v are highlighted.
\begin{tikzpicture}[every node/.style={font=\scriptsize}]
  \foreach \i in {0,...,15} { \node[circle,fill=black,inner sep=0pt,minimum size=2.5pt] (m\i) at (180+360/16*\i:1.6) {}; }
  \draw[cC,line width=.25pt,opacity=.55] (m0) -- (m1);
  \draw[cB,line width=.25pt,opacity=.55] (m0) -- (m2);
  \draw[cD,line width=.25pt,opacity=.55] (m0) -- (m4);
  \draw[cA,line width=.25pt,opacity=.55] (m0) -- (m5);
  \draw[cA,line width=.25pt,opacity=.55] (m0) -- (m6);
  \draw[cA,line width=.25pt,opacity=.55] (m0) -- (m7);
  \draw[cC,line width=.25pt,opacity=.55] (m0) -- (m8);
  \draw[cA,line width=.25pt,opacity=.55] (m0) -- (m9);
  \draw[cB,line width=.25pt,opacity=.55] (m0) -- (m10);
  \draw[cA,line width=.25pt,opacity=.55] (m0) -- (m12);
  \draw[cA,line width=.25pt,opacity=.55] (m0) -- (m13);
  \draw[cD,line width=.25pt,opacity=.55] (m0) -- (m14);
  \draw[cD,line width=.25pt,opacity=.55] (m0) -- (m15);
  \draw[cA,line width=.25pt,opacity=.55] (m1) -- (m2);
  \draw[cB,line width=.25pt,opacity=.55] (m1) -- (m3);
  \draw[cA,line width=.25pt,opacity=.55] (m1) -- (m4);
  \draw[cD,line width=.25pt,opacity=.55] (m1) -- (m5);
  \draw[cA,line width=.25pt,opacity=.55] (m1) -- (m6);
  \draw[cA,line width=.25pt,opacity=.55] (m1) -- (m7);
  \draw[cB,line width=.25pt,opacity=.55] (m1) -- (m8);
  \draw[cA,line width=.25pt,opacity=.55] (m1) -- (m9);
  \draw[cD,line width=.25pt,opacity=.55] (m1) -- (m10);
  \draw[cA,line width=.25pt,opacity=.55] (m1) -- (m11);
  \draw[cB,line width=.25pt,opacity=.55] (m1) -- (m12);
  \draw[cA,line width=.25pt,opacity=.55] (m1) -- (m13);
  \draw[cB,line width=.25pt,opacity=.55] (m1) -- (m14);
  \draw[cC,line width=.25pt,opacity=.55] (m1) -- (m15);
  \draw[cD,line width=.25pt,opacity=.55] (m2) -- (m3);
  \draw[cB,line width=.25pt,opacity=.55] (m2) -- (m4);
  \draw[cA,line width=.25pt,opacity=.55] (m2) -- (m5);
  \draw[cC,line width=.25pt,opacity=.55] (m2) -- (m6);
  \draw[cB,line width=.25pt,opacity=.55] (m2) -- (m7);
  \draw[cA,line width=.25pt,opacity=.55] (m2) -- (m8);
  \draw[cB,line width=.25pt,opacity=.55] (m2) -- (m9);
  \draw[cC,line width=.25pt,opacity=.55] (m2) -- (m10);
  \draw[cA,line width=.25pt,opacity=.55] (m2) -- (m11);
  \draw[cA,line width=.25pt,opacity=.55] (m2) -- (m12);
  \draw[cA,line width=.25pt,opacity=.55] (m2) -- (m13);
  \draw[cB,line width=.25pt,opacity=.55] (m2) -- (m14);
  \draw[cD,line width=.25pt,opacity=.55] (m2) -- (m15);
  \draw[cD,line width=.25pt,opacity=.55] (m3) -- (m4);
  \draw[cC,line width=.25pt,opacity=.55] (m3) -- (m5);
  \draw[cD,line width=.25pt,opacity=.55] (m3) -- (m6);
  \draw[cD,line width=.25pt,opacity=.55] (m3) -- (m7);
  \draw[cC,line width=.25pt,opacity=.55] (m3) -- (m9);
  \draw[cC,line width=.25pt,opacity=.55] (m3) -- (m10);
  \draw[cB,line width=.25pt,opacity=.55] (m3) -- (m11);
  \draw[cB,line width=.25pt,opacity=.55] (m3) -- (m12);
  \draw[cB,line width=.25pt,opacity=.55] (m3) -- (m13);
  \draw[cA,line width=.25pt,opacity=.55] (m3) -- (m14);
  \draw[cC,line width=.25pt,opacity=.55] (m3) -- (m15);
  \draw[cD,line width=.25pt,opacity=.55] (m4) -- (m5);
  \draw[cC,line width=.25pt,opacity=.55] (m4) -- (m6);
  \draw[cD,line width=.25pt,opacity=.55] (m4) -- (m7);
  \draw[cC,line width=.25pt,opacity=.55] (m4) -- (m8);
  \draw[cA,line width=.25pt,opacity=.55] (m4) -- (m9);
  \draw[cA,line width=.25pt,opacity=.55] (m4) -- (m10);
  \draw[cD,line width=.25pt,opacity=.55] (m4) -- (m11);
  \draw[cB,line width=.25pt,opacity=.55] (m4) -- (m12);
  \draw[cC,line width=.25pt,opacity=.55] (m4) -- (m13);
  \draw[cB,line width=.25pt,opacity=.55] (m4) -- (m14);
  \draw[cD,line width=.25pt,opacity=.55] (m4) -- (m15);
  \draw[cD,line width=.25pt,opacity=.55] (m5) -- (m6);
  \draw[cA,line width=.25pt,opacity=.55] (m5) -- (m7);
  \draw[cA,line width=.25pt,opacity=.55] (m5) -- (m8);
  \draw[cC,line width=.25pt,opacity=.55] (m5) -- (m9);
  \draw[cC,line width=.25pt,opacity=.55] (m5) -- (m10);
  \draw[cC,line width=.25pt,opacity=.55] (m5) -- (m11);
  \draw[cD,line width=.25pt,opacity=.55] (m5) -- (m12);
  \draw[cD,line width=.25pt,opacity=.55] (m5) -- (m13);
  \draw[cA,line width=.25pt,opacity=.55] (m5) -- (m14);
  \draw[cA,line width=.25pt,opacity=.55] (m5) -- (m15);
  \draw[cC,line width=.25pt,opacity=.55] (m6) -- (m7);
  \draw[cD,line width=.25pt,opacity=.55] (m6) -- (m8);
  \draw[cA,line width=.25pt,opacity=.55] (m6) -- (m9);
  \draw[cA,line width=.25pt,opacity=.55] (m6) -- (m10);
  \draw[cC,line width=.25pt,opacity=.55] (m6) -- (m11);
  \draw[cD,line width=.25pt,opacity=.55] (m6) -- (m12);
  \draw[cC,line width=.25pt,opacity=.55] (m6) -- (m13);
  \draw[cD,line width=.25pt,opacity=.55] (m6) -- (m14);
  \draw[cC,line width=.25pt,opacity=.55] (m6) -- (m15);
  \draw[cA,line width=.25pt,opacity=.55] (m7) -- (m8);
  \draw[cD,line width=.25pt,opacity=.55] (m7) -- (m9);
  \draw[cC,line width=.25pt,opacity=.55] (m7) -- (m10);
  \draw[cB,line width=.25pt,opacity=.55] (m7) -- (m11);
  \draw[cA,line width=.25pt,opacity=.55] (m7) -- (m12);
  \draw[cD,line width=.25pt,opacity=.55] (m7) -- (m13);
  \draw[cA,line width=.25pt,opacity=.55] (m7) -- (m14);
  \draw[cB,line width=.25pt,opacity=.55] (m7) -- (m15);
  \draw[cC,line width=.25pt,opacity=.55] (m8) -- (m9);
  \draw[cB,line width=.25pt,opacity=.55] (m8) -- (m10);
  \draw[cB,line width=.25pt,opacity=.55] (m8) -- (m12);
  \draw[cD,line width=.25pt,opacity=.55] (m8) -- (m13);
  \draw[cD,line width=.25pt,opacity=.55] (m8) -- (m14);
  \draw[cD,line width=.25pt,opacity=.55] (m8) -- (m15);
  \draw[cA,line width=.25pt,opacity=.55] (m9) -- (m10);
  \draw[cB,line width=.25pt,opacity=.55] (m9) -- (m11);
  \draw[cD,line width=.25pt,opacity=.55] (m9) -- (m12);
  \draw[cD,line width=.25pt,opacity=.55] (m9) -- (m13);
  \draw[cC,line width=.25pt,opacity=.55] (m9) -- (m14);
  \draw[cB,line width=.25pt,opacity=.55] (m9) -- (m15);
  \draw[cD,line width=.25pt,opacity=.55] (m10) -- (m11);
  \draw[cC,line width=.25pt,opacity=.55] (m10) -- (m12);
  \draw[cD,line width=.25pt,opacity=.55] (m10) -- (m13);
  \draw[cC,line width=.25pt,opacity=.55] (m10) -- (m14);
  \draw[cD,line width=.25pt,opacity=.55] (m10) -- (m15);
  \draw[cB,line width=.25pt,opacity=.55] (m11) -- (m12);
  \draw[cB,line width=.25pt,opacity=.55] (m11) -- (m13);
  \draw[cA,line width=.25pt,opacity=.55] (m11) -- (m14);
  \draw[cB,line width=.25pt,opacity=.55] (m11) -- (m15);
  \draw[cB,line width=.25pt,opacity=.55] (m12) -- (m13);
  \draw[cB,line width=.25pt,opacity=.55] (m12) -- (m14);
  \draw[cB,line width=.25pt,opacity=.55] (m12) -- (m15);
  \draw[cA,line width=.25pt,opacity=.55] (m13) -- (m14);
  \draw[cD,line width=.25pt,opacity=.55] (m13) -- (m15);
  \draw[cB,line width=.25pt,opacity=.55] (m14) -- (m15);
  \draw[cA,line width=1.3pt] (m0) -- (m11);
  \draw[cB,line width=1.3pt] (m8) -- (m11);
  \draw[cA,line width=1.3pt] (m0) -- (m3);
  \draw[cB,line width=1.3pt] (m3) -- (m8);
  \node[left=0pt] at (m0) {$u$};
  \node[right=0pt] at (m8) {$v$};
  \node[above right=-1pt] at (m11) {$z$};
  \node[below left=-1pt] at (m3) {$z'$};
\end{tikzpicture}


    \NOTE{$\M_0$: randomly colored, $O(k\,4^k)$ vertices \fcite{jmt1995}}
  \end{column}
  \begin{column}{.60\textwidth}
    \KEY{Two-witness property} \ For all $u \ne v$ and all colors $\textcolor{cA}{\boldsymbol{\mu}}, \textcolor{cB}{\boldsymbol{\nu}}$,
    there are at least two \emph{witnesses} $z$ such that
    the $u$-$z$ and $z$-$v$ edges have colors $\textcolor{cA}{\boldsymbol{\mu}}$ and $\textcolor{cB}{\boldsymbol{\nu}}$, respectively.
  \end{column}
\end{columns}

\vfill
\pause
\vspace{.3em}
\KEY{Every \CoR{} term $t$ is color-invariant on $\M_0$}: $\M_0 \models t(u, v) \iff \M_0 \models t(u', v')$
whenever the $u$-$v$ and $u'$-$v'$ edges have the same color.

\vspace{1em}
Proof by induction on $t$: the case, $t = t_1 \mathbin{;} t_2$, follows from the (two-)witness property.

\note{
\textbf{スライド概要}：$\M_0$ ではどの色の道のパターンにも中間の頂点がある。だから項は「色 $\mapsto$ T/F」の表になる。

\textbf{メモ}：
\begin{itemize}
  \item 例：$P_1 \mathbin{;} Q_1$ はどの辺でも T（色 $\set{1}$ と $\emptyset$ の道が必ずある）、ループでは F
  \item 「2 つ以上」が効くのは $\M_\lambda$ 側（hub を避ける中間の頂点が残る）
\end{itemize}
}
\end{frame}

% ============================================================
\begin{frame}{$\M_\lambda$: recoloring at the hub $\hub$}

% Figure 1 of the paper, redrawn around the mechanism: the read edge at the hub
% accepts only witness edges of color lambda, and the fold recolors exactly those.
% Colors as on the slide of G_k: read edge cA, witness edges cB.
\begin{center}
\begin{tikzpicture}[every node/.style={font=\small},x=1.1cm,y=1.1cm,
    xmark/.style={text=red,font=\large\bfseries,inner sep=0pt}]
  \foreach \sx/\name/\lab/\vis in {0/0/structure $\M_0$/1-,7.3/1/recolored structure $\M_\lambda$/2-} {
    \begin{scope}[xshift=\sx cm,visible on=<\vis>]
      \node at (.6,2.35) {\lab};
      \node[pv,label=below:$\hub$] (h\name) at (0,0) {};
      \node[pv,label=below:$v$] (v\name) at (2.4,0) {};
      \node[pv] (z\name) at (-.5,1.6) {};
      \node[pv] (w\name) at (.6,1.75) {};
      \node[pv] (o\name) at (-1.6,.3) {};
      \draw[black!35] (h\name) -- node[below,font=\scriptsize,text=black!50] {$\mu \ne \lambda$} (o\name);
      \draw[line width=1.4pt,cA] (h\name) -- node[below=2pt,font=\scriptsize] {color $\lambda^c$} (v\name);
    \end{scope}
  }
  % M_0: witness edges of color lambda at the hub
  \draw[thick,cB] (h0) -- node[left,font=\scriptsize] {$\lambda$} (z0);
  \draw[thick,cB] (h0) -- node[right,font=\scriptsize] {$\lambda$} (w0);
  \node[font=\scriptsize,text=cB] at (-.2,2.0) {color $\lambda$};
  \node at (1.2,-1.05) {$G_k(\hub, v) = 1$};
  % M_lambda: the same edges, recolored to lambda^c, are no longer witnesses
  \begin{scope}[visible on=<2->]
    \draw[thick,black!35,densely dashed] (h1) -- node[left,font=\scriptsize,text=black!60] {$\lambda^c$} (z1);
    \draw[thick,black!35,densely dashed] (h1) -- node[right,font=\scriptsize,text=black!60] {$\lambda^c$} (w1);
    \node[xmark] at ($(h1)!.75!(z1)$) {$\times$};
    \node[xmark] at ($(h1)!.75!(w1)$) {$\times$};
    \node at (7.3/1.1+1.2,-1.05) {$G_k(\hub, v) = \textcolor{red}{0}$};
    \draw[-{Stealth},thick] (3.3,1.0) -- node[above] {recolor $\lambda$} (4.6,1.0);
  \end{scope}
\end{tikzpicture}
\end{center}

\vfill
\begin{center}
  $G_k(\hub, v)$ needs a $\hub$-$z$ edge of color $\lambda$
  \onslide<2->{\\ recoloring $\lambda$ changes \emph{exactly those edges at $\hub$}, \ and nothing else}
\end{center}

\vfill
\onslide<3->{%
\begin{center}
  $2^k$ recolored structures $\M_\lambda$, one for each color $\lambda$ \ \NOTE{($=$ the $\M_1, \dots, \M_{2^k}$ of the outline)}
  \\ away from $\hub$, terms cannot distinguish $\M_\lambda$ from $\M_0$ \NOTE{(of the $\ge 2$ middle vertices, one avoids $\hub$)}
\end{center}}

\note{
\textbf{スライド概要}：塗り替えで $G_k$ が 1 から 0 に変わる。

\textbf{メモ}：
\begin{itemize}
  \item $\hub$ に接する色 $\lambda$ の辺だけ塗り替える
  \item 色ごとに 1 つ、計 $2^k$ 個
\end{itemize}
}
\end{frame}

% ============================================================
\begin{frame}{Why a composition newly distinguishes at most two $\M_\lambda$}

% Figure 2 of the paper, redrawn with one layout for all panels:
% the pair is the bottom edge, witnesses sit above (and below) it.
% Red with a cross: an edge recolored by the fold; thick green: a witness that survives.
\begin{center}
\begin{tikzpicture}[every node/.style={font=\small},x=1.2cm,y=1.2cm,
    xmark/.style={text=red,font=\large\bfseries,inner sep=0pt},
    lost/.style={thick,red,densely dashed},
    kept/.style={line width=1.4pt,cC}]
  % (a) M_0
  \begin{scope}
    \node at (.9,1.75) {(a) $\M_0$, at $(\hub, v)$};
    \node[pv,label=below:$\hub$] (h) at (0,0) {};
    \node[pv,label=above:$z$] (z) at (.9,1.0) {};
    \node[pv,label=below:$v$] (v) at (1.8,0) {};
    \draw[thick] (h) -- node[above left] {$\lambda$} (z);
    \draw[thick] (z) -- node[above right] {$\nu$} (v);
    \draw[densely dotted] (h) -- (v);
    \node[align=center] at (.9,-.85) {middle vertex $z$:\\ $(X \mathbin{;} Y)(\hub, v) = 1$};
  \end{scope}
  % (b) M_lambda, at the hub: every witness uses an edge at the hub
  \begin{scope}[xshift=4.4cm,visible on=<2->]
    \node at (.9,1.75) {(b) $\M_\lambda$, at $(\hub, v)$};
    \node[pv,label=below:$\hub$] (h) at (0,0) {};
    \node[pv,label=above:$z$] (z) at (.9,1.0) {};
    \node[pv,label=below:$v$] (v) at (1.8,0) {};
    \draw[lost] (h) -- node[above left] {$\lambda^c$} (z);
    \node[xmark] at ($(h)!.5!(z)$) {$\times$};
    \draw[thick] (z) -- node[above right] {$\nu$} (v);
    \draw[densely dotted] (h) -- (v);
    \node[align=center] at (.9,-.85) {every middle vertex is lost:\\ \alert{distinguished}};
  \end{scope}
  % (c) M_lambda, away from the hub: the hub is only one of the witnesses
  \begin{scope}[xshift=8.8cm,visible on=<3->]
    \node at (.9,1.75) {(c) $\M_\lambda$, at $(u, v)$};
    \node[pv,label=left:$u$] (u) at (0,0) {};
    \node[pv,label=above:$\hub$] (h) at (.9,1.0) {};
    \node[pv,label=right:$v$] (v) at (1.8,0) {};
    \node[pv,label=below:$z'$] (z) at (.9,-.8) {};
    \draw[lost] (u) -- node[above left] {$\lambda^c$} (h);
    \node[xmark] at ($(u)!.5!(h)$) {$\times$};
    \draw[thick] (h) -- node[above right] {$\nu$} (v);
    \draw[kept] (u) -- node[below left] {$\lambda$} (z);
    \draw[kept] (z) -- node[below right] {$\nu$} (v);
    \draw[densely dotted] (u) -- (v);
    \node[align=center] at (.9,-1.45) {a second middle vertex $z'$\\ avoids $\hub$: \ not distinguished};
  \end{scope}
\end{tikzpicture}
\end{center}
\NOTE{\hfill a composition $X \mathbin{;} Y$, \ where $X$ holds on the edges of color $\lambda$ and $Y$ on those of color $\nu \ne \lambda$}

\vfill
\onslide<4->{%
\begin{center}
  only pairs at the hub can change, and only if $X$ or $Y$ singles out the color $\lambda$
  \\ $\Rightarrow$ \ each composition newly distinguishes at most \alert{two} $\M_\lambda$ from $\M_0$ \ \NOTE{(hence the Key Lemma)}
\end{center}}

\note{
\textbf{スライド概要}：$\hub$ を含む点対では見分けられうる、含まない点対では見分けられない。

\textbf{メモ}：
\begin{itemize}
  \item $\hub$ を含まない点対では、$\hub$ を通らない中間の頂点 $z'$ が残る
  \item 条件は 3 つ、どれも $\lambda$ を決める $\to$ 高々 2
\end{itemize}
}
\end{frame}
